\documentclass[11pt]{amsart}
\usepackage{amsmath,amssymb,amsthm,mathtools}
\usepackage[margin=2.7cm]{geometry}
\usepackage[colorlinks=true,linkcolor=blue,citecolor=blue]{hyperref}
\newtheorem{theorem}{Theorem}
\newtheorem{lemma}[theorem]{Lemma}
\newtheorem{remark}[theorem]{Remark}
\title{Fourier decay of the Gelfand--Leray measure on an odd superellipsoid}
\author{}\date{}
\begin{document}\maketitle

Let \(d\ge2\), let \(k\ge3\) be odd, and set
\[
P(x)=\sum_{j=1}^d|x_j|^k,\qquad
\Sigma_k=\{x\in\mathbb R^d:P(x)=1\}.
\]
The Gelfand--Leray measure \(\nu\) is defined by
\begin{equation}\label{eq:GL}
dP\wedge d\nu=dx_1\wedge\cdots\wedge dx_d,\qquad
d\nu=\frac{d\mathcal H^{d-1}}{|\nabla P|}.
\end{equation}
Since \(\nabla P\ne0\) on \(\Sigma_k\), this is a finite positive measure.
We use
\[
\widehat\nu(\xi)=\int_{\Sigma_k}e^{-i\langle\xi,x\rangle}\,d\nu(x).
\]

\begin{theorem}\label{thm:main}
There is \(C=C(d,k)\) such that
\begin{equation}\label{eq:main}
|\widehat\nu(\xi)|\le C(1+|\xi|)^{-\alpha},\qquad
\alpha=\min\left\{k,\frac{d-1}{k}\right\}.
\end{equation}
Equivalently, for \(|\xi|\ge1\),
\[
|\widehat\nu(\xi)|
\le C\max\bigl\{|\xi|^{-k},|\xi|^{-(d-1)/k}\bigr\}.
\]
In particular, if \(d\le k^2+1\), then
\[
|\widehat\nu(\xi)|\le C(1+|\xi|)^{-(d-1)/k}.
\]
\end{theorem}

The proof uses the following standard local form of the convex finite-type
Fourier estimate. It is stated explicitly so that a radial one-dimensional
argument is not mistaken for a multidimensional estimate.

\begin{lemma}[piecewise finite-type estimate]\label{lem:finite-type}
Let \(m\ge1\), \(Q(y)=\sum_{j=1}^m|y_j|^k\), and let
\(g\in C^\infty((-\delta,\delta))\) satisfy \(|g'|\ge c>0\).
Suppose that \(a\) is compactly supported in \(\{Q<\delta\}\), is smooth on
each open orthant, has bounded one-sided derivatives through order \(k\),
and has matching one-sided jets through order \(k-1\) at every coordinate
hyperplane. Then, uniformly for \(s\in\mathbb R^m\) and \(R\ge1\),
\begin{equation}\label{eq:FT-lemma}
\left|\int_{\mathbb R^m}e^{iR(g(Q(y))+s\cdot y)}a(y)\,dy\right|
\le C_a\bigl(R^{-m/k}+R^{-k}\bigr).
\end{equation}
\end{lemma}

\begin{proof}
On each orthant, \(Q\) is polynomial and \(g\circ Q\) has finite line type
\(k\). A tangent plane cuts off a cap of height \(R^{-1}\) whose
\(m\)-dimensional measure is \(O(R^{-m/k})\), uniformly in the tangent
plane: at the maximally flat point this is the dilation
\(y=R^{-1/k}z\), and elsewhere convexity gives a smaller cap. The local
cap theorem of Bruna--Nagel--Wainger \cite[Theorems A and B]{BNW} gives
\(O(R^{-m/k})\) for the smooth interior of each orthant, uniformly after
addition of the linear function \(s\cdot y\).

To reassemble the orthants, truncate them at distance \(\varepsilon\) from
their faces and use the integration-by-parts construction in
\cite[Section~6]{BNW}. On adding adjacent orthants, face terms of orders
below \(k\) cancel because the one-sided jets agree. The first possible
uncancelled term and the corresponding remainder are \(O(R^{-k})\).
The same argument is applied successively at intersections of faces.
The bounds are independent of \(\varepsilon\); letting
\(\varepsilon\downarrow0\) proves \eqref{eq:FT-lemma}.
\end{proof}

\begin{proof}[Proof of Theorem~\ref{thm:main}]
For \(|\xi|\le1\), the assertion follows from \(\nu(\Sigma_k)<\infty\).
Assume \(R=|\xi|\ge1\) and write \(\xi=Rv\), \(v\in S^{d-1}\).

Choose \(c_d>0\) so small that the sets
\[
U_{j,\sigma}=\{x\in\Sigma_k:\sigma x_j>c_d\},
\qquad 1\le j\le d,\quad \sigma\in\{\pm1\},
\]
cover \(\Sigma_k\), and choose a subordinate partition of unity with the
same piecewise-\(C^k\) regularity as \(P\).
It suffices to treat \(U_{1,+}\). Put
\[
y=(x_2,\ldots,x_d),\quad
Q(y)=\sum_{j=2}^d|y_j|^k,\quad F(y)=(1-Q(y))^{1/k}.
\]
Thus \(x_1=F(y)\), and \eqref{eq:GL} gives the exact chart formula
\begin{equation}\label{eq:GL-chart}
d\nu=\frac{dy}{kF(y)^{k-1}}.
\end{equation}
In particular, the Euclidean graph Jacobian
\(\sqrt{1+|\nabla F|^2}\) must not be used here.

A localized contribution is
\begin{equation}\label{eq:local}
J(R,v)=\int e^{-iR(v_1F(y)+v'\cdot y)}a(y)\,dy,
\end{equation}
where \(a\) contains the partition function and
\((kF^{k-1})^{-1}\).

Suppose first that \(|v_1|\ge c_0>0\). After replacing \(R\) by
\(R|v_1|\), the phase is
\[
g(Q(y))+s\cdot y,\quad
g(u)=\operatorname{sgn}(v_1)(1-u)^{1/k},\quad s=v'/|v_1|.
\]
On the compact chart, \(|g'|\) is bounded below. The amplitude is smooth
on every orthant, has bounded one-sided derivatives through order \(k\),
and has matching jets through order \(k-1\), since every nonsmooth
occurrence of \(y_j\) is through \(|y_j|^k\).
Lemma~\ref{lem:finite-type}, with \(m=d-1\), gives
\[
|J(R,v)|\le C\bigl(R^{-(d-1)/k}+R^{-k}\bigr)
\le2C R^{-\min\{k,(d-1)/k\}}.
\]

If \(|v_1|<c_0\), choose \(c_0\) below the positive lower bound for the
first component of the unit normal on the support of this chart. The phase
in \eqref{eq:local} then has no critical point, and compactness gives
\[
|\nabla_y(v_1F(y)+v'\cdot y)|\ge c>0.
\]
Integrating by parts \(k\) times on truncated orthants, using
\[
L=\frac{\nabla_y(v_1F+v'\cdot y)}
{iR|\nabla_y(v_1F+v'\cdot y)|^2}\cdot\nabla_y,
\]
gives \(O(R^{-k})\). Boundary terms below order \(k\) cancel between
adjacent orthants because the one-sided jets match; hence the truncation
may tend to zero.

The same argument applies to every \(U_{j,\sigma}\). The cover is finite
and the constants are uniform in \(v\), so summing proves
\eqref{eq:main}.
\end{proof}

\begin{remark}
For odd \(k\), \(|t|^k\) is \(C^{k-1}\), not \(C^k\), at zero. Thus one
may neither apply a globally smooth stationary-phase theorem nor integrate
by parts across \(t=0\) without tracking one-sided boundary terms.
Moreover, one-dimensional van der Corput in a radial variable gives only
\(R^{-1/k}\); by itself it does not give \(R^{-(d-1)/k}\).
\end{remark}

\begin{thebibliography}{9}
\bibitem{BNW}
J.~Bruna, A.~Nagel, and S.~Wainger,
\emph{Convex hypersurfaces and Fourier transforms},
Ann. of Math. (2) \textbf{127} (1988), no.~2, 333--365.
\end{thebibliography}
\end{document}
